
\section*{Abstract} Maritime safety and surveillance has always been a critical challenge, particularly when vessels lose communication due to equipment failure, harsh weather conditions, or intentional signal suppression. However current monitoring systems have several limitations that could not be solved properly in the past research. A primary research gap is the inability to detect small ships, specifically those under 20 meters. Additionally existing researches struggle with sea clutter, the noise caused by waves and rough weather leading to frequent missed detection or false alarm. Furthermore, most approaches fail to maintain temporal consistency. They cannot reliably track a vessel’s movement over time to reconstruct its navigation history. Without this missing trajectory information, it becomes difficult to predict a missing ship's location or what went wrong during an incident. However, in our research we plan to explore the possible ways to overcome these gaps, by introducing a multimodal ship detection framework that works with radar signals and satellite imageries to improve the maritime safety operations all together. By combining information from SAR data and satellite images we will explore the possibility of accurately pinpointing ships of any size by analyzing multiple images and identifying actual ships that have consistent appearances throughout all the images from both sources, allowing to reduce false positives from sea clutter. Additionally we intend on exploring further into the idea of storing the sequential coordinates of any neighbouring vessels’ to create a trail which will allow us to observe and analyze their movement. In our experiment, we have selected five dual-stream RGB-SAR detectors to train and examine the performance of our selected MOS-ship dataset~\cite{you2026multi}. Among all the five models, which are: Oriented R-CNN~\cite{xie2021oriented}, S2A-net~\cite{han2021align}, OrientedFormer~\cite{zhao2024orientedformer}, YOLO26-OBB~\cite{jocher2026ultralytics} and Faster R-CNN~\cite{7485869}, the model YOLO26-OBB performed the best with an mAP50:95 of 0.871 while being the smallest (33.2 M parameters) and fastest (about 25 FPS).



\vspace{1cm}
\textbf{Keywords:} Maritime surveillance, ship detection, small vessel detection, synthetic aperture radar (SAR), optical satellite imagery, multimodal fusion, sea clutter, false alarm reduction, temporal consistency, vessel tracking, trajectory reconstruction, missing ship prediction, satellite image analysis, deep learning-based detection.
\pagebreak
